\documentclass{article}
\usepackage{amsmath, amsthm, amssymb}
\usepackage{hyperref}
\usepackage{cleveref}

\newtheorem{theorem}{Theorem}
\newtheorem{lemma}[theorem]{Lemma}
\newtheorem{definition}[theorem]{Definition}
\newtheorem{corollary}[theorem]{Corollary}
\newtheorem{remark}[theorem]{Remark}

\title{AI Assurance as a Category-Theoretic Preservation Structure:\\
       A Lean 4 Formalization}
\author{GhostDrift Mathematical Institute}
\date{2026}

\begin{document}

\maketitle

\section{Introduction}

AI governance in multi-organizational deployments is difficult because operational claims often cross institutional boundaries. A model may be trained by one organization, integrated by another, monitored by a third, and audited by a fourth. When governance claims rest only on self-attestation, there is no structural guarantee that the evidence needed for inspection remains available after these systems are composed.

The main claim formalized here is that inspectable AI governance requires a faithful evidence layer. Without such a layer, governance-relevant distinctions, including different judgment grounds, responsibility paths, and audit traces, can collapse to the same operational appearance. This is a structural necessity, not a policy preference.

This paper presents a Lean 4 formalization of this claim. We construct the evidence category as the Grothendieck construction of an indexed fiber family over the operational category, prove the relevant preservation theorems, and exhibit a finite counterexample establishing the necessity direction.

\section{Mathematical Preliminaries}

A category consists of objects, morphisms between objects, identity morphisms, and associative composition. A functor maps objects and morphisms from one category to another while preserving identities and composition. A functor is faithful when its action on each hom-set is injective.

Given a functor $F : O \to \mathbf{Cat}$, the Grothendieck construction $\int_O F$ is the category whose objects are pairs $(X,a)$ with $X \in O$ and $a \in F_X$, and whose morphisms $(X,a) \to (Y,b)$ are pairs $(f,\alpha)$ with $f : X \to Y$ and $\alpha : F(f)(a) \to b$.

An adjunction $L \dashv R$ between two categories is a natural bijection between morphisms $L X \to Y$ and morphisms $X \to RY$. A section of a forgetful functor $U : E \to O$ is a functor $S : O \to E$ such that $U \circ S$ is the identity functor on $O$, using the composition convention of the Lean formalization.

\begin{lemma}[faithful\_of\_section]
If $S : O \to E$ and $U : E \to O$ satisfy $U \circ S = \mathrm{Id}_O$, then $S$ is faithful.
\end{lemma}

\section{The AI Assurance Structure}

\begin{definition}[IndexedAssurance]
An \emph{IndexedAssurance} structure is a tuple
\[
  (O,F,\mathrm{push},\mathrm{pull},\dashv,
   \mathrm{push\_id},\mathrm{push\_comp},
   \mathrm{pull\_id},\mathrm{pull\_comp},
   s,\mathrm{standard\_push},BC)
\]
with the following components:
\begin{itemize}
\item $O$ is the operational category.
\item $F_X$ is the evidence fiber category at $X \in O$.
\item $f_!$ is the push functor for $f : X \to Y$.
\item $f^*$ is the pull functor for $f : X \to Y$.
\item $f_! \dashv f^*$ is an adjunction.
\item $\mathrm{push\_id}$ states that $f_! = \mathrm{Id}$ when $f = \mathrm{id}_X$.
\item $\mathrm{push\_comp}$ states that $(g \circ f)_! = g_! \circ f_!$.
\item $\mathrm{pull\_id}$ states that $f^* = \mathrm{Id}$ when $f = \mathrm{id}_X$.
\item $\mathrm{pull\_comp}$ states that $(g \circ f)^* = f^* \circ g^*$.
\item $s_X$ is the standard evidence at $X$.
\item $\mathrm{standard\_push}$ states that $f_!(s_X)=s_Y$.
\item $BC$ is the Beck-Chevalley condition for allowed squares, retained as a structural field.
\end{itemize}
\end{definition}

\section{The Derived Evidence Category}

The evidence category is derived as
\[
  E = \int_O F,
\]
the Grothendieck construction of the Lean definition \texttt{toCatFunctor}. Objects are pairs $(X,a)$ with $X \in O$ and $a \in F_X$. Morphisms are pairs $(f,\alpha)$ with $f : X \to Y$ and $\alpha : f_!(a) \to b$.

\begin{definition}[standardSection]
The standard section $S : O \to E$ is defined by $S(X)=(X,s_X)$ and $S(f)=(f,\mathrm{id}_{s_Y})$, where the identity is transported along \texttt{standard\_push}.
\end{definition}

\begin{theorem}[section\_eq]
The Lean formalization proves
\[
  U \circ S = \mathrm{Id}_O.
\]
\end{theorem}

\begin{corollary}[section\_faithful]
The standard section $S$ is faithful.
\end{corollary}

\section{Universal Properties}

\begin{theorem}[opcart\_factor]
For any $h : (X,a) \to (Z,c)$ in $E$ with $\mathrm{base}(h)=f \circ g$, there exists a unique
\[
  \delta : (Y,f_!(a)) \to (Z,c)
\]
with $\mathrm{base}(\delta)=g$ and
\[
  \mathrm{opcartLift}(f,a) \circ \delta = h.
\]
\end{theorem}

\begin{theorem}[cart\_factor]
For any $h : (Z,c) \to (Y,b)$ in $E$ with $\mathrm{base}(h)=k \circ f$, there exists a unique
\[
  \delta : (Z,c) \to (X,f^*(b))
\]
with $\mathrm{base}(\delta)=k$ and
\[
  \delta \circ \mathrm{cartLift}(f,b) = h.
\]
\end{theorem}

The theorem \texttt{opcart\_factor} means that any evidence path over a composed operation factors uniquely through the forward record update. The theorem \texttt{cart\_factor} means that any audit path factors uniquely through the backward audit pullback. Together they state that evidence propagation and audit retrieval are universal, not ad hoc.

\section{The Necessity Counterexample}

\begin{definition}[CollapseCounterexample]
Let $E$ be the category with objects $\{\mathrm{src},\mathrm{tgt}\}$ and morphisms $\{\mathrm{idSrc},\mathrm{idTgt},\mathrm{traceA},\mathrm{traceB}\}$, where $\mathrm{traceA}\neq\mathrm{traceB}$. Let $O$ be the category with objects $\{\mathrm{src},\mathrm{tgt}\}$ and morphisms $\{\mathrm{idSrc},\mathrm{idTgt},\mathrm{op}\}$. Let $U : E \to O$ map both $\mathrm{traceA}$ and $\mathrm{traceB}$ to $\mathrm{op}$.
\end{definition}

\begin{theorem}[trace\_distinction\_collapses]
\[
  \mathrm{traceA}\neq \mathrm{traceB}
  \quad\text{and}\quad
  U(\mathrm{traceA}) = U(\mathrm{traceB}).
\]
\end{theorem}

\begin{theorem}[U\_not\_faithful]
The functor $U$ is not faithful.
\end{theorem}

\begin{corollary}[forgetting\_trace\_layer\_can\_collapse\_distinctions]
There exist categories $E$, $O$ and a functor $U : E \to O$ such that $U$ is not faithful.
\end{corollary}

This counterexample shows that operational sameness does not imply governance sameness. Two morphisms representing distinct judgment grounds or responsibility paths can be identified by a forgetful functor. Therefore, if AI governance requires those distinctions to remain inspectable after composition, a faithful evidence layer is structurally necessary.

\section{Limitations and Future Work}

\begin{itemize}
\item \texttt{beck\_chevalley} is retained as a structural field. It is not derived from \texttt{adj + push\_comp + pull\_comp} in general. A full proof requires additional coherence conditions specific to each instance.
\item The formalization covers the mathematical core only. It does not address the translation from formal model to operational practice, real-world AI deployments, or regulatory compliance.
\item Future work: replace \texttt{beck\_chevalley} with a derived proof for specific instances, for example presheaf categories and polynomial functors. Extend to lax natural transformations for non-strict split conditions.
\end{itemize}

\section{Conclusion}

The formalization proves that the Grothendieck construction gives a derived evidence category rather than a freely postulated one. The standard section is faithful, opcartesian and cartesian universal properties hold at the TotalCategory level, and the collapse counterexample establishes the necessity direction.

This formalization establishes that AI assurance, understood as the preservation of evidence, responsibility paths, and judgment grounds across organizational composition, has a precise mathematical characterization that is machine-checkable. The claim is not that any deployed AI system satisfies this structure, but that any governance framework claiming inspectability must instantiate it.

\begin{thebibliography}{9}

\bibitem{maclane} Mac Lane, S. \emph{Categories for the Working Mathematician}. Springer, 1971.
\bibitem{grothendieck} Grothendieck, A. \emph{Revêtements étales et groupe fondamental (SGA 1)}. 1960.
\bibitem{mathlib} The Mathlib Community. Mathlib4. \url{https://github.com/leanprover-community/mathlib4}
\bibitem{euaiact} European Parliament. Regulation (EU) 2024/1689 (AI Act). 2024.
\bibitem{nist} NIST. \emph{AI Risk Management Framework 1.0}. 2023.
\bibitem{jacobs} Jacobs, B. \emph{Categorical Logic and Type Theory}. Elsevier, 1999.

\end{thebibliography}

\end{document}
